\section{Conclusion}
\label{sec:conclusion}

\subsection{Summary of Findings}

This paper has shown that T1 managed freight lanes on the Interstate 2.0 network create the physical precondition for intercity bus service that does not currently exist in most of the continental United States: scheduled, reliable, and affordable transit at 62 mph average speed with Planning Time Index $\leq$ 1.15.

Across 12 T1 corridors enabled by the F.1 hub network \citep{ROUTE_F1}:
\begin{itemize}
  \item T1 bus travel times are \textbf{28--45\% shorter} than current bus alternatives on every corridor where current bus service exists, with a simple average improvement of 33\%.
  \item T1 bus is \textbf{competitive with or faster than Amtrak} on all corridors where Amtrak service operates, and provides service on 7 corridors where Amtrak offers no direct routing.
  \item \textbf{Five corridors} (Atlanta--Dallas, Houston--Chicago, Atlanta--Jacksonville interior, Denver--Salt Lake City, Miami--Atlanta interior) gain entirely new intercity connectivity that no existing mode provides.
  \item At a targeted fare of \$0.12/mile, T1 express bus is the \textbf{lowest-cost intercity travel option} in every market it serves, while being dramatically more reliable than the incumbent Greyhound service it would replace.
  \item The 12 corridors together project \textbf{24 million annual passengers} at market equilibrium, supporting 38 bus operators and 12,000 direct jobs.
\end{itemize}

\subsection{The Reliability Insight}

The core finding of this paper is methodological as much as empirical: PTI is not merely a freight metric. It is the infrastructure precondition for scheduled transit. The transition from PTI 1.86 (current I-80) to PTI 1.15 (I2.0 target) is not a 38\% improvement in congestion --- it is the difference between a market that operators rationally exit and a market they compete to enter.

This has a policy implication that extends beyond the corridors analyzed here: any managed lane investment that achieves PTI $\leq$ 1.20 on a corridor connecting populations of 500,000 or more at each end is, implicitly, a transit investment, regardless of whether bus service is explicitly planned. The infrastructure creates the market; the market will find the operators.

\subsection{Policy Recommendation: Bus Access Provisions}

The clearest policy lever is the concession agreement. When T1 managed lane segments are built or franchised, concession agreements should require guaranteed bus operator access at a per-mile access fee set to cover incremental lane capacity costs only --- not to generate profit for the concessionaire. This is analogous to the open-access provisions required of railroad companies under the Interstate Commerce Act. A managed lane that excludes scheduled bus operators is infrastructure built on public right-of-way that systematically advantages private automobile and truck use over transit --- a policy choice that should be explicit, not implicit.

For marginal corridors, a designation analogous to Essential Air Service (49 U.S.C. § 41731), which would require new authorization under the Surface Transportation Reauthorization Act, provides a mechanism for public subsidy at a fraction of the cost of Essential Air Service: \$15/passenger versus \$175/passenger, on corridors where the communities served have higher C3 scores and lower vehicle ownership than the typical EAS community \citep{EAS_DOT2023, ACS_B08201_2022}. The closest existing authority is 49 U.S.C. § 5311(f) (FTA rural intercity bus program), which funds rural intercity bus connections but does not mandate service frequency or cap fares.

\subsection{The I2.0 Transit Obligation}

Taken together, F.1 \citep{ROUTE_F1} and F.2 make the following case: the Interstate 2.0 transit layer is not an optional add-on to a freight investment. It is an obligation created by the investment itself. When \$253 billion of public funds are committed to upgrading the highway system, and that upgrade creates the physical precondition for intercity transit to 12.4 million transit-dependent Americans \citep{ROUTE_F1} and 24 million annual passengers on 12 corridors, excluding bus access is a deliberate policy choice to leave those Americans out of the public investment. That choice requires justification that does not exist.

\subsection{Limitations}

The 24 million annual passenger estimate uses a gravity model calibrated on existing intercity bus routes \citep{BTS_intermodal2023}. New markets --- corridors with no current service --- have no calibration data; we applied a 30\% new-market discount, but this is an assumption rather than a measurement. First-year ridership on genuinely new corridors (Atlanta--Dallas, Houston--Chicago) may be substantially below the equilibrium projection. The Essential Intercity Transportation subsidy analysis should be treated as order-of-magnitude; actual subsidy requirements will depend on first-year ridership outcomes and operator cost structures specific to each corridor.

A companion paper (F.3, planned) will analyze multimodal integration at T1/T1 hub stations: the design standard for bus-to-freight-truck transfer, intercity bus-to-local transit connections, and the capital cost of full intermodal facility development.
